\documentclass[12pt,a4paper]{article}
\usepackage[utf8]{inputenc}
\usepackage[T1]{fontenc}
\usepackage[italian]{babel}
\usepackage{geometry}
\usepackage{longtable}
\usepackage{array}
\usepackage{booktabs}
\usepackage{hyperref}
\usepackage{xcolor}
\usepackage{listings}
\usepackage{enumitem}

\geometry{margin=2.2cm}

\lstdefinestyle{sqlstyle}{
  basicstyle=\ttfamily\small,
  keywordstyle=\color{blue},
  commentstyle=\color{gray},
  stringstyle=\color{teal},
  showstringspaces=false,
  breaklines=true,
  frame=single
}

\title{Documentazione Completa e Didattica del Progetto\\Concessionario Auto \& Moto}
\author{N86005495 \and N86005712}
\date{\today}

\begin{document}

\maketitle
\tableofcontents
\newpage

\section{Come leggere questo documento (se parti da zero)}
Questo documento e stato scritto per una persona che non conosce Java.
Ogni elemento ha sempre tre parti:
\begin{itemize}
  \item \textbf{Cosa fa}: descrizione tecnica sintetica.
  \item \textbf{Perche esiste}: motivazione pratica della presenza nel progetto.
  \item \textbf{Spiegazione semplice}: spiegazione senza dare per scontati concetti avanzati.
\end{itemize}

\section{Fondamenti minimi (Java, Swing, JDBC) in linguaggio semplice}
\subsection{Che cos'e Java (in due righe)}
Java e un linguaggio di programmazione. Un programma Java e fatto da classi e metodi.
Il metodo speciale \texttt{main} e il punto di partenza: quando avvii il programma, Java entra li.

\subsection{Che cos'e Swing}
Swing e una libreria grafica di Java per creare finestre, pulsanti, caselle di testo e tabelle.
Nel progetto, Swing serve per l'interfaccia del concessionario.

\subsection{Che cos'e JDBC}
JDBC e il modo standard con cui Java parla con il database SQL.
In pratica: Java apre una connessione, invia query SQL, riceve risultati.

\section{Panoramica architetturale}
\textbf{Cosa fa:} il progetto implementa un gestionale desktop per concessionario auto/moto con login a ruoli.

\textbf{Perche esiste:} separare ruoli, dati e schermate rende il programma ordinato, testabile e piu facile da mantenere.

\textbf{Spiegazione semplice:} immagina tre piani:
\begin{enumerate}
  \item Piano interfaccia: cio che l'utente vede (finestre e pulsanti).
  \item Piano logica: regole applicative (chi puo fare cosa).
  \item Piano dati: database PostgreSQL dove i dati vengono salvati.
\end{enumerate}

Flusso generale:
\begin{enumerate}
  \item Avvio da \texttt{Main.main()}.
  \item Delega a \texttt{MainApp.main()}.
  \item Creazione finestra principale \texttt{ConcessionarioFrame}.
  \item Login per ruolo (cliente o dipendente).
  \item Navigazione con \texttt{CardLayout} + \texttt{JTabbedPane}.
  \item Chiamate ai DAO e query SQL.
  \item Vincoli logici garantiti anche lato DB (trigger/funzioni).
\end{enumerate}

\section{Classi Java e metodi con motivazione e spiegazione}

\subsection{Package it.unina.prog}

\subsubsection*{DBManager}
\textbf{Cosa fa:}
\begin{itemize}
  \item \texttt{getConnection(): Connection} apre una connessione JDBC al database PostgreSQL.
\end{itemize}

\textbf{Perche esiste:}
serve un punto unico per aprire connessioni al DB, cosi si evita di duplicare URL, utente e password in molti file.

\textbf{Spiegazione semplice:}
DBManager e come una porta di ingresso al database: ogni volta che un DAO deve leggere/scrivere dati, passa da qui.

\textbf{Chiamata principale usata e motivo:}
\begin{itemize}
  \item \texttt{DriverManager.getConnection(url, user, password)}: e la chiamata standard Java per collegarsi al DB.
\end{itemize}

\subsubsection*{Main}
\textbf{Cosa fa:}
\begin{itemize}
  \item \texttt{main(String[] args): void} richiama \texttt{MainApp.main(args)}.
\end{itemize}

\textbf{Perche esiste:}
fornisce un punto di ingresso pulito e semplice; puo essere utile anche per cambiare bootstrap senza toccare il resto.

\textbf{Spiegazione semplice:}
Main e il campanello di casa: suoni li e poi vieni reindirizzato nella stanza dove parte davvero l'app.

\subsubsection*{MainApp}
\textbf{Cosa fa:}
\begin{itemize}
  \item Avvia la GUI sul thread corretto di Swing.
  \item Imposta il look and feel del sistema.
  \item Mostra \texttt{ConcessionarioFrame}.
\end{itemize}

\textbf{Perche esiste:}
in Swing l'interfaccia deve essere costruita su un thread specifico (EDT), altrimenti possono comparire bug grafici.

\textbf{Spiegazione semplice:}
MainApp prepara il palco prima dello spettacolo: decide stile grafico e poi apre la finestra principale.

\textbf{Chiamate principali e perche:}
\begin{itemize}
  \item \texttt{SwingUtilities.invokeLater(...)}: garantisce che la GUI parta sul thread corretto.
  \item \texttt{UIManager.setLookAndFeel(...)}: usa stile nativo (piu familiare all'utente).
  \item \texttt{new ConcessionarioFrame().setVisible(true)}: crea e mostra la finestra.
\end{itemize}

\subsection{Package it.unina.prog.model}

\subsubsection*{Cliente}
\textbf{Cosa fa:}
contiene i dati di un cliente (id, nome, tipo, email, telefono).

\textbf{Perche esiste:}
serve un contenitore chiaro per passare dati tra DAO e interfaccia.

\textbf{Spiegazione semplice:}
e una scheda anagrafica in memoria: non fa calcoli complessi, conserva solo informazioni.

\textbf{Metodi e motivazione:}
\begin{itemize}[leftmargin=1.2cm]
  \item \texttt{Cliente(...)}: costruttore. Perche: crea l'oggetto completo in un colpo solo.
  \item \texttt{getId()}, \texttt{getNome()}, \texttt{getTipo()}, \texttt{getEmail()}, \texttt{getTelefono()}: getter. Perche: permettono lettura sicura dei dati senza esporre i campi direttamente.
\end{itemize}

\subsubsection*{Dipendente}
\textbf{Cosa fa:}
modella un dipendente con id, nome, ruolo e password.

\textbf{Perche esiste:}
serve per autenticare l'utente interno e per mostrare ruolo/id nella sessione.

\textbf{Spiegazione semplice:}
e il ``cartellino'' digitale del dipendente che entra nel sistema.

\textbf{Metodi e motivazione:}
\begin{itemize}
  \item Costruttore \texttt{Dipendente(...)}: crea il record in memoria.
  \item Getter \texttt{getId/getNome/getRuolo/getPassword}: usati da login e interfaccia.
\end{itemize}

\subsubsection*{Veicolo}
\textbf{Cosa fa:}
rappresenta un veicolo con dati anagrafici, prezzo e stato.

\textbf{Perche esiste:}
serve una struttura unica per riempire tabelle, combo box e operazioni di vendita/test drive.

\textbf{Spiegazione semplice:}
e la ``carta d'identita'' del mezzo dentro l'app.

\textbf{Metodi e motivazione:}
\begin{itemize}
  \item Costruttore \texttt{Veicolo(...)}: crea un veicolo completo.
  \item Getter \texttt{getTarga/getMarca/getModello/getTipo/getPrezzo/getStato}: rendono disponibili i dati a UI e DAO.
\end{itemize}

\subsection{Package it.unina.prog.dao}

\subsubsection*{ClienteDAO}
\textbf{Cosa fa:}
gestisce le operazioni SQL legate ai clienti.

\textbf{Perche esiste:}
separare SQL dalla GUI evita codice disordinato e facilita manutenzione.

\textbf{Spiegazione semplice:}
DAO significa Data Access Object: e il ``commesso'' che va in magazzino (DB), prende o modifica dati e li riporta all'app.

\textbf{Metodi, perche e spiegazione:}
\begin{itemize}
  \item \texttt{inserisciCliente(...)}: crea un cliente nel DB. Perche: serve al pannello dipendente quando registra un nuovo cliente.
  \item \texttt{getClienti()}: ritorna lista clienti. Perche: serve per mostrare tabella e combo selezione.
  \item \texttt{getClienteByEmail(email)}: cerca cliente per email. Perche: utile in validazioni/lookup.
  \item \texttt{autenticaCliente(email, password)}: verifica credenziali cliente. Perche: login ruolo cliente.
  \item \texttt{aggiornaCliente(...)}: modifica dati senza cambiare password. Perche: aggiornamenti ordinari.
  \item \texttt{aggiornaClienteConPassword(...)}: modifica anche password. Perche: gestione completa profilo.
  \item \texttt{eliminaCliente(id)}: elimina cliente. Perche: operazione amministrativa.
\end{itemize}

\textbf{Chiamate JDBC ricorrenti e motivo:}
\begin{itemize}
  \item \texttt{prepareStatement}: evita SQL injection e gestisce parametri.
  \item \texttt{setString/setInt/...}: inserisce valori nelle query.
  \item \texttt{executeQuery/executeUpdate}: esegue SELECT o INSERT/UPDATE/DELETE.
\end{itemize}

\subsubsection*{DipendenteDAO}
\textbf{Cosa fa:}
recupera e autentica dipendenti.

\textbf{Perche esiste:}
la logica di accesso dipendente non deve stare nella GUI.

\textbf{Spiegazione semplice:}
quando un dipendente inserisce id+password, questa classe controlla se esiste davvero.

\textbf{Metodi:}
\begin{itemize}
  \item \texttt{getDipendenteById(id)}: legge record da DB.
  \item \texttt{autenticaDipendente(id, password)}: confronta password e ritorna oggetto valido o \texttt{null}.
\end{itemize}

\subsubsection*{VeicoloDAO}
\textbf{Cosa fa:}
contiene CRUD veicoli e calcolo stato.

\textbf{Perche esiste:}
il catalogo veicoli e centrale: vendite, test drive e vista catalogo dipendono da questi dati.

\textbf{Spiegazione semplice:}
e il reparto ``parco auto/moto'': inserisce mezzi, li aggiorna, li elimina e dice se sono disponibili.

\textbf{Metodi con perche:}
\begin{itemize}
  \item \texttt{inserisciVeicolo(...)}: crea nuovo veicolo.
  \item \texttt{getStatoVeicolo(targa)}: chiede al DB lo stato corrente. Perche: lo stato dipende da vendite/test drive.
  \item \texttt{getVeicoliDisponibili()}: filtra solo disponibili. Perche: in vendita o test drive non puoi proporre mezzi non disponibili.
  \item \texttt{aggiornaVeicolo(...)}: aggiorna anagrafica e prezzo.
  \item \texttt{eliminaVeicolo(targa)}: rimuove mezzo.
  \item \texttt{getAllVeicoli()}: elenco completo per catalogo e gestione.
\end{itemize}

\subsubsection*{TestDriveDAO}
\textbf{Cosa fa:}
inserisce richieste test drive.

\textbf{Perche esiste:}
la prenotazione test drive e una transazione dedicata.

\textbf{Spiegazione semplice:}
quando scegli cliente, veicolo e data, questa classe salva la prenotazione.

\textbf{Metodo:}
\begin{itemize}
  \item \texttt{richiediTestDrive(clienteId, targa, data)}.
\end{itemize}

\subsubsection*{VenditaDAO}
\textbf{Cosa fa:}
registra una vendita in database.

\textbf{Perche esiste:}
la vendita e evento critico: va tracciata con data, cliente, dipendente, prezzo.

\textbf{Spiegazione semplice:}
quando confermi la vendita in interfaccia, questa classe scrive la riga ufficiale nel DB.

\textbf{Metodo:}
\begin{itemize}
  \item \texttt{effettuaVendita(targa, clienteId, dipendenteId, prezzoFinale)}.
\end{itemize}

\subsection{Package it.unina.prog.ui.common}

\subsubsection*{UiSupport}
\textbf{Cosa fa:}
fornisce funzioni comuni per stile pannelli e gestione errori utente.

\textbf{Perche esiste:}
evita copia-incolla di codice grafico e messaggi d'errore in molti file.

\textbf{Spiegazione semplice:}
e la ``cassetta attrezzi'' condivisa della GUI.

\textbf{Metodi con motivazione:}
\begin{itemize}
  \item \texttt{wrapTitled(panel, title)}: mette bordo con titolo uniforme. Perche: coerenza visiva.
  \item \texttt{showErr(parent, ex)}: mostra errore standard. Perche: centralizzare UX degli errori.
  \item \texttt{toUserMessage(ex)}: traduce errori tecnici in frasi comprensibili. Perche: utente non deve leggere stack trace.
\end{itemize}

\subsection{Package it.unina.prog.ui.model}

\subsubsection*{UiModels}
\textbf{Cosa fa:}
contiene classi ``item'' per mostrare valori utili in combo box e liste.

\textbf{Perche esiste:}
se metti solo id o targa grezzi in una combo, l'esperienza utente peggiora.

\textbf{Spiegazione semplice:}
queste classi preparano etichette leggibili tipo ``Mario Rossi (3)'' o ``AB123CD - Yaris''.

\textbf{Elementi:}
\begin{itemize}
  \item \texttt{ClienteItem}: tiene id e nome; \texttt{toString()} definisce testo mostrato.
  \item \texttt{VeicoloItem}: tiene targa, modello e prezzo; \texttt{toString()} definisce testo mostrato.
  \item \texttt{ClienteLogin}: contenitore cliente+password (supporto flussi login).
\end{itemize}

\subsection{Package it.unina.prog.ui.validation}

\subsubsection*{InputValidator}
\textbf{Cosa fa:}
raccolta di validazioni comuni su email, telefono, targa, date.

\textbf{Perche esiste:}
la validazione ripetuta in piu pannelli genera bug e incoerenza; centralizzarla evita errori.

\textbf{Spiegazione semplice:}
invece di riscrivere gli stessi controlli ovunque, c'e una sola classe che li fa bene una volta.

\textbf{Metodi e motivazione:}
\begin{itemize}
  \item \texttt{isValidEmail(email)}: impedisce email palesemente malformate.
  \item \texttt{isValidPhone(phone)}: forza formato coerente (10 cifre).
  \item \texttt{normalizePlate(plate)}: uniforma targa (trim+maiuscolo).
  \item \texttt{isValidPlate(plate)}: blocca targhe con caratteri/spazi non ammessi.
  \item \texttt{parseDateOrThrow(dateInput)}: converte stringa in data o lancia errore chiaro.
  \item \texttt{requireTodayOrFuture(date)}: vieta date passate.
\end{itemize}

\subsection{Package it.unina.prog.ui}

\subsubsection*{ConcessionarioFrame}
\textbf{Cosa fa:}
e la finestra principale e il coordinatore di tutta la navigazione.

\textbf{Perche esiste:}
serve un ``orchestratore'' che sappia quale schermata mostrare in base al login e allo stato sessione.

\textbf{Spiegazione semplice:}
pensa a un centralino: riceve accesso utente e lo instrada nella sezione giusta.

\textbf{Metodi con motivazione e spiegazione:}
\begin{itemize}
  \item \texttt{ConcessionarioFrame()}: inizializza struttura base. Perche: impostare finestra una sola volta.
  \item \texttt{createHeader()}: crea barra titolo/sessione/logout. Perche: info utente sempre visibili.
  \item \texttt{createLoginPanel()}: costruisce form login e listener. Perche: ingresso unico app.
  \item \texttt{buildDipendenteArea(id)}: genera tab operative dipendente. Perche: separare strumenti interni.
  \item \texttt{buildClienteArea(cliente)}: genera tab cliente. Perche: limitare azioni ai permessi cliente.
  \item \texttt{switchToLogin()}: torna alla schermata iniziale.
  \item \texttt{switchToDipendente(dipendente)}: monta vista dipendente e aggiorna sessione.
  \item \texttt{switchToCliente(cliente)}: monta vista cliente e aggiorna sessione.
  \item \texttt{handleClienteLogin(email,password)}: valida input e autentica cliente.
  \item \texttt{handleDipendenteLogin(idText,password)}: valida input e autentica dipendente.
\end{itemize}

\section{Pannelli cliente con spiegazione didattica}

\subsection{CatalogoClientePanel}
\textbf{Cosa fa:} mostra l'elenco dei veicoli in tabella.

\textbf{Perche esiste:} il cliente deve poter consultare catalogo e disponibilita senza modificare dati.

\textbf{Spiegazione semplice:} e una vetrina in sola lettura.

\textbf{Chiamate importanti e motivo:}
\begin{itemize}
  \item \texttt{VeicoloDAO.getAllVeicoli()}: recupera dati aggiornati.
  \item \texttt{DefaultTableModel}: struttura righe/colonne tabella.
  \item \texttt{Runnable load}: permette refresh rapido senza duplicare codice.
\end{itemize}

\subsection{ProfiloClientePanel}
\textbf{Cosa fa:} mostra e aggiorna email/telefono del cliente loggato.

\textbf{Perche esiste:} il cliente deve poter mantenere aggiornati i propri contatti.

\textbf{Spiegazione semplice:} e la sezione ``i miei dati''.

\textbf{Elementi chiave e motivo:}
\begin{itemize}
  \item \texttt{validateEmailTelefono()}: evita salvataggi con dati invalidi.
  \item SELECT profilo per riempire campi all'apertura.
  \item UPDATE profilo al click su salva.
\end{itemize}

\subsection{TestDriveClientePanel}
\textbf{Cosa fa:} permette prenotazione test drive lato cliente.

\textbf{Perche esiste:} il test drive e azione principale prima dell'acquisto.

\textbf{Spiegazione semplice:} il cliente sceglie veicolo+data e invia richiesta.

\textbf{Chiamate chiave e motivo:}
\begin{itemize}
  \item \texttt{VeicoloDAO.getVeicoliDisponibili()}: solo mezzi prenotabili.
  \item \texttt{InputValidator.parseDateOrThrow()} + \texttt{requireTodayOrFuture()}: evita errori data.
  \item \texttt{TestDriveDAO.richiediTestDrive()}: salvataggio prenotazione.
\end{itemize}

\subsection{OperazioniClientePanel}
\textbf{Cosa fa:} visualizza storico vendite e test drive del cliente.

\textbf{Perche esiste:} trasparenza: il cliente deve vedere le proprie operazioni.

\textbf{Spiegazione semplice:} e l'estratto conto delle attivita cliente.

\textbf{Perche query dirette qui:}
si tratta di viste aggregate con join e ordinamenti specifici, comode per report UI immediato.

\section{Pannelli dipendente con spiegazione didattica}

\subsection{ClientiPanel}
\textbf{Cosa fa:} CRUD completo sui clienti.

\textbf{Perche esiste:} i dipendenti devono registrare, correggere o rimuovere anagrafiche.

\textbf{Spiegazione semplice:} e il ``gestionale clienti''.

\textbf{Metodi/chiamate e motivo:}
\begin{itemize}
  \item \texttt{validateClienteInput(...)}: controlla campi obbligatori e coerenza.
  \item \texttt{ClienteDAO.inserisci/aggiorna/elimina/getClienti}: operazioni base archivio clienti.
\end{itemize}

\subsection{VeicoliPanel}
\textbf{Cosa fa:} CRUD veicoli + gestione marca non presente.

\textbf{Perche esiste:} inventario veicoli e cuore del concessionario.

\textbf{Spiegazione semplice:} e il ``magazzino'' auto/moto.

\textbf{Metodi/chiamate e motivo:}
\begin{itemize}
  \item \texttt{validateVeicoloInput()}: blocca inserimenti sporchi.
  \item \texttt{targaExistsIgnoringCaseAndSpaces()}: evita duplicati furbi (spazi/case).
  \item \texttt{ensureMarcaExistsOrCreate()}: permette aggiunta marca al volo.
  \item \texttt{VeicoloDAO.*}: CRUD veicoli.
\end{itemize}

\subsection{VenditaPanel}
\textbf{Cosa fa:} registra vendita con dipendente corrente.

\textbf{Perche esiste:} trasformare disponibilita in vendita registrata in modo tracciabile.

\textbf{Spiegazione semplice:} e la cassa finale: scegli cliente e veicolo, confermi, la vendita entra nello storico.

\textbf{Chiamate chiave:}
\begin{itemize}
  \item \texttt{VeicoloDAO.getVeicoliDisponibili()} e \texttt{ClienteDAO.getClienti()} per popolare combo.
  \item \texttt{VenditaDAO.effettuaVendita(...)} per salvataggio definitivo.
\end{itemize}

\subsection{TestDrivePanel}
\textbf{Cosa fa:} inserisce test drive da parte dipendente.

\textbf{Perche esiste:} anche il personale interno deve poter registrare prenotazioni assistite.

\textbf{Spiegazione semplice:} versione operatore della prenotazione test drive.

\textbf{Chiamate:}
\begin{itemize}
  \item stesse del pannello cliente (veicoli disponibili, clienti, validazione data, insert test drive).
\end{itemize}

\subsection{OrdiniPanel}
\textbf{Cosa fa:} mostra storico globale (vendite + test drive) e permette eliminazioni.

\textbf{Perche esiste:} controllo amministrativo e correzione dati errati.

\textbf{Spiegazione semplice:} e il ``registro generale'' del concessionario.

\textbf{Chiamate e motivo:}
\begin{itemize}
  \item query JOIN per avere dati leggibili con nomi (non solo id).
  \item DELETE per correggere record selezionati.
\end{itemize}

\section{Chiamate esterne ricorrenti (con motivazione)}

\subsection{Swing/AWT}
\textbf{Perche usate:} costruiscono interfaccia desktop nativa Java.

\textbf{Spiegazione semplice:}
\begin{itemize}
  \item \texttt{JPanel}: contenitore di componenti.
  \item \texttt{Layout} (Border/GridBag/Card/Flow): decide come disporre gli elementi.
  \item \texttt{JButton/JTextField/JTable}: input e visualizzazione dati.
  \item \texttt{JOptionPane}: popup messaggi/confirm.
  \item \texttt{addActionListener}: reagisce al click utente.
\end{itemize}

\subsection{JDBC}
\textbf{Perche usato:} e il ponte standard tra Java e SQL.

\textbf{Spiegazione semplice:}
\begin{enumerate}
  \item apri connessione,
  \item prepari query con parametri,
  \item esegui,
  \item leggi risultato,
  \item chiudi risorse (automatico con try-with-resources).
\end{enumerate}

\subsection{Utility Java base}
\textbf{Perche usate:}
\begin{itemize}
  \item parsing numeri/date: convertire testo input in tipi utili.
  \item metodi stringa: pulizia e validazione input utente.
\end{itemize}

\section{Query SQL dirette nei panel: cosa fanno e perche}

\subsection{ProfiloClientePanel}
\textbf{Perche:} leggere e aggiornare profilo del cliente loggato.
\begin{lstlisting}[style=sqlstyle,language=SQL]
SELECT nome, tipo, email, telefono
FROM Cliente
WHERE id = ?;

UPDATE Cliente
SET email = ?, telefono = ?
WHERE id = ?;
\end{lstlisting}

\subsection{VeicoliPanel}
\textbf{Perche:} controllo duplicati targa, verifica marca, eventuale inserimento marca.
\begin{lstlisting}[style=sqlstyle,language=SQL]
SELECT 1
FROM Veicolo
WHERE UPPER(REPLACE(targa, ' ', '')) = UPPER(REPLACE(?, ' ', ''));

SELECT 1
FROM Marca
WHERE LOWER(nome) = LOWER(?);

INSERT INTO Marca (nome, nazione)
VALUES (?, ?);
\end{lstlisting}

\subsection{OperazioniClientePanel}
\textbf{Perche:} mostrare storico cliente con dati completi e ordinati.
\begin{lstlisting}[style=sqlstyle,language=SQL]
SELECT v.id, v.data, v.veicolo, v.prezzo_finale, d.nome AS dipendente
FROM Vendita v
JOIN Dipendente d ON v.dipendente = d.id
WHERE v.cliente = ?
ORDER BY v.data DESC;

SELECT id, data, veicolo
FROM TestDrive
WHERE cliente = ?
ORDER BY data DESC;
\end{lstlisting}

\subsection{OrdiniPanel}
\textbf{Perche:} vista amministrativa globale + cancellazioni mirate.
\begin{lstlisting}[style=sqlstyle,language=SQL]
SELECT v.id, v.data, v.veicolo, c.nome, d.nome, v.prezzo_finale
FROM Vendita v
JOIN Cliente c ON v.cliente = c.id
JOIN Dipendente d ON v.dipendente = d.id
ORDER BY v.data DESC;

SELECT t.id, t.data, c.nome, t.veicolo
FROM TestDrive t
JOIN Cliente c ON t.cliente = c.id
ORDER BY t.data DESC;

DELETE FROM Vendita WHERE id = ?;
DELETE FROM TestDrive WHERE id = ?;
\end{lstlisting}

\section{Funzioni SQL, trigger e vincoli con motivazione}

\subsection{Funzione get\_stato\_veicolo(targa)}
\textbf{Cosa fa:}
ritorna \texttt{venduto}, \texttt{non disponibile momentaneamente} o \texttt{disponibile}.

\textbf{Perche esiste:}
lo stato del mezzo dipende da regole dati (vendite/test drive), quindi e corretto calcolarlo nel DB.

\textbf{Spiegazione semplice:}
il database decide in autonomia lo stato reale del veicolo, cosi la UI non inventa regole diverse.

\subsection{Trigger check\_veicolo\_disponibile su Vendita}
\textbf{Cosa fa:} blocca la seconda vendita dello stesso veicolo.

\textbf{Perche esiste:} evita incoerenza gravissima (veicolo venduto due volte).

\textbf{Spiegazione semplice:} e un ``controllore automatico'' dentro il DB prima dell'inserimento.

\subsection{Trigger check\_testdrive\_limite su TestDrive}
\textbf{Cosa fa:} blocca oltre 3 test drive mensili per cliente.

\textbf{Perche esiste:} regola di business imposta dal dominio.

\textbf{Spiegazione semplice:} il DB conta prenotazioni del mese e stoppa la quarta.

\subsection{Vincoli principali schema}
\textbf{Perche esistono:}
per impedire dati impossibili o incoerenti direttamente alla fonte.

\textbf{Esempi spiegati:}
\begin{itemize}
  \item CHECK su tipo cliente/veicolo: evita valori fuori elenco.
  \item UNIQUE su email cliente: evita duplicati identita.
  \item FOREIGN KEY: collega tabelle in modo coerente.
  \item ON DELETE CASCADE: elimina automaticamente dati figli collegati.
\end{itemize}

\section{Mappa completa delle chiamate (con perche)}
\begin{enumerate}
  \item \textbf{Avvio:} \texttt{Main.main()} $\rightarrow$ \texttt{MainApp.main()} $\rightarrow$ \texttt{ConcessionarioFrame()}.
  Perche: separa ingresso tecnico e avvio GUI.

  \item \textbf{Login cliente:} \texttt{handleClienteLogin()} $\rightarrow$ \texttt{ClienteDAO.autenticaCliente()} $\rightarrow$ \texttt{switchToCliente()}.
  Perche: validazione input, verifica credenziali, apertura area corretta.

  \item \textbf{Login dipendente:} \texttt{handleDipendenteLogin()} $\rightarrow$ \texttt{DipendenteDAO.autenticaDipendente()} $\rightarrow$ \texttt{switchToDipendente()}.
  Perche: stesso principio del cliente ma con ruolo interno.

  \item \textbf{Tab cliente} (profilo, catalogo, test drive, storico).
  Perche: separa operazioni in blocchi semplici da capire.

  \item \textbf{Tab dipendente} (clienti, veicoli, vendita, test drive, ordini).
  Perche: strumenti operativi completi solo per personale autorizzato.
\end{enumerate}

\section{Perche sono stati deduplicati i doppioni}
\textbf{Scelta:} pattern identici (esempio: struttura JDBC, listener Swing, refresh con \texttt{Runnable load}) sono spiegati una sola volta.

\textbf{Perche:}
\begin{itemize}
  \item il documento resta studiabile,
  \item si evita ripetizione inutile,
  \item si capisce meglio la logica comune del progetto.
\end{itemize}

\textbf{Spiegazione semplice:}
se dieci classi fanno la stessa cosa nello stesso modo, capire bene una volta basta per capirle tutte.

\section{Glossario minimo per principianti}
\begin{itemize}
  \item \textbf{Classe}: modello/ricetta di un oggetto.
  \item \textbf{Metodo}: azione che una classe puo eseguire.
  \item \textbf{Oggetto}: istanza concreta di una classe.
  \item \textbf{DAO}: classe che accede al database.
  \item \textbf{Query}: istruzione SQL (SELECT, INSERT, UPDATE, DELETE).
  \item \textbf{Trigger}: regola automatica del DB che scatta su eventi.
  \item \textbf{GUI}: interfaccia grafica.
  \item \textbf{CRUD}: Create, Read, Update, Delete.
\end{itemize}

\section{Conclusione}
Questa versione aggiunge per ogni sezione e per ogni elemento principale:
\begin{itemize}
  \item spiegazione tecnica,
  \item motivazione del perche esiste,
  \item spiegazione didattica per chi non conosce Java.
\end{itemize}

In questo modo il file puo essere usato sia come documentazione tecnica sia come materiale di studio introduttivo.

\newpage
\section{Appendice molto approfondita: studio guidato file per file}
Questa appendice rende la documentazione molto piu lunga e molto piu dettagliata.
L'obiettivo non e solo dire \emph{cosa} fa il codice, ma anche spiegare \emph{come ragiona} e \emph{perche} ogni pezzo esiste.

\subsection{1. File di avvio}

\subsubsection*{Main.java}
\textbf{Ruolo nel progetto:} e il punto di ingresso piu esterno, cioe il primo file che la macchina cerca quando parte il programma.

\textbf{Perche serve davvero:} in molti progetti si usa una classe molto piccola come questa per avere un avvio semplice e uniforme. Se in futuro si volesse cambiare la vera classe di avvio, basterebbe modificare un solo punto.

\textbf{Spiegazione per principianti:} immagina una porta d'ingresso di un edificio. Quando il programma parte, entra da qui e poi viene accompagnato verso l'app vera e propria.

\textbf{Metodo presente:}
\begin{itemize}
  \item \texttt{main(String[] args)}: legge gli argomenti di avvio e li passa a \texttt{MainApp}.
\end{itemize}

\textbf{Motivazione della chiamata:} chiamare \texttt{MainApp.main(args)} rende esplicito che la logica di avvio reale non e duplicata, ma centralizzata.

\subsubsection*{MainApp.java}
\textbf{Ruolo nel progetto:} gestisce l'avvio grafico vero e proprio.

\textbf{Perche serve davvero:} Swing richiede attenzione al thread dell'interfaccia. Se si apre la finestra nel modo sbagliato, l'app puo diventare instabile o comportarsi male.

\textbf{Spiegazione per principianti:} questo file e come il direttore di scena: prepara tutto prima che la finestra sia visibile.

\textbf{Passaggi interni spiegati uno per uno:}
\begin{enumerate}
  \item \texttt{SwingUtilities.invokeLater(...)}: dice a Java di costruire la finestra nel thread corretto.
  \item \texttt{UIManager.setLookAndFeel(...)}: cambia l'aspetto dei pulsanti e delle finestre per renderlo simile al sistema operativo.
  \item \texttt{new ConcessionarioFrame()}: crea la finestra principale.
  \item \texttt{setVisible(true)}: rende la finestra visibile sullo schermo.
\end{enumerate}

\textbf{Perche queste chiamate sono importanti:} senza di esse la UI non partirebbe nel modo piu corretto e la prima impressione del programma sarebbe peggiore.

\subsection{2. File di collegamento al database}

\subsubsection*{DBManager.java}
\textbf{Ruolo nel progetto:} e il punto unico in cui si apre la connessione al database.

\textbf{Perche serve davvero:} il database e un servizio esterno; invece di ripetere la stessa configurazione in ogni classe, si centralizza tutto qui.

\textbf{Spiegazione per principianti:} pensa a DBManager come al telefono del database. Ogni volta che il programma vuole chiedere qualcosa, deve prima fare questo numero.

\textbf{Metodo presente:}
\begin{itemize}
  \item \texttt{getConnection()}: restituisce un oggetto \texttt{Connection} pronto per essere usato.
\end{itemize}

\textbf{Motivazione della scelta tecnica:}
un metodo unico rende piu semplice modificare host, porta, nome database o credenziali senza cercare in tanti file diversi.

\textbf{Cosa bisogna capire a livello base:}
\begin{itemize}
  \item \texttt{Connection} e la sessione attiva con il database.
  \item Se la connessione e aperta, si possono inviare query SQL.
  \item Se la connessione fallisce, il resto del programma non puo leggere o scrivere dati.
\end{itemize}

\subsection{3. Modelli dati}

\subsubsection*{Cliente.java}
\textbf{Ruolo nel progetto:} rappresenta il cliente come oggetto Java.

\textbf{Perche serve davvero:} i dati letti dal database devono avere una forma leggibile dal codice. Un oggetto Cliente e piu utile di tante variabili sparse.

\textbf{Spiegazione per principianti:} e come una scheda cartacea di un cliente, ma in versione digitale.

\textbf{Costruttore e getter spiegati:}
\begin{itemize}
  \item \texttt{Cliente(...)}: crea il cliente con tutti i dati gia pronti.
  \item \texttt{getId()}: restituisce l'identificativo numerico.
  \item \texttt{getNome()}: restituisce il nome.
  \item \texttt{getTipo()}: dice se il cliente e privato o azienda.
  \item \texttt{getEmail()}: restituisce la mail.
  \item \texttt{getTelefono()}: restituisce il numero di telefono.
\end{itemize}

\textbf{Perche i getter esistono:} servono a leggere i dati senza modificare direttamente i campi interni.

\subsubsection*{Dipendente.java}
\textbf{Ruolo nel progetto:} rappresenta il dipendente che puo entrare nel gestionale.

\textbf{Perche serve davvero:} il login dipendente deve ottenere un oggetto che descrive l'utente autenticato, non solo un numero o una stringa.

\textbf{Spiegazione per principianti:} e il badge del lavoratore: contiene chi e, che ruolo ha e quale password usa per accedere.

\textbf{Metodi principali:}
\begin{itemize}
  \item Costruttore \texttt{Dipendente(...)}: crea il dipendente completo.
  \item Getter dei campi: permettono alla GUI di mostrare nome, ruolo e id.
\end{itemize}

\textbf{Perche la password e nel modello:} il programma la usa nel login. In un progetto didattico questo e semplice da capire; in un progetto reale sarebbe meglio salvarla in modo sicuro.

\subsubsection*{Veicolo.java}
\textbf{Ruolo nel progetto:} rappresenta un mezzo disponibile nel catalogo.

\textbf{Perche serve davvero:} il catalogo, le vendite e i test drive lavorano tutti sugli stessi dati veicolo.

\textbf{Spiegazione per principianti:} e la scheda del veicolo: targa, marca, modello, tipo, prezzo e stato.

\textbf{Metodi principali:}
\begin{itemize}
  \item Costruttore \texttt{Veicolo(...)}: crea il mezzo con tutti i dettagli.
  \item Getter: restituiscono i singoli campi quando servono in tabelle o combo box.
\end{itemize}

\subsection{4. DAO: il ponte con il database}

\subsubsection*{ClienteDAO.java}
\textbf{Ruolo nel progetto:} si occupa di tutte le operazioni SQL sui clienti.

\textbf{Perche serve davvero:} senza un DAO, la SQL finirebbe dentro la grafica, rendendo il codice confuso e difficile da correggere.

\textbf{Spiegazione per principianti:} e il lavoratore che va a parlare con il database al posto dell'interfaccia grafica.

\paragraph{inserisciCliente(nome, tipo, email, telefono, password)}
\textbf{Cosa fa:} salva un nuovo cliente nel database.

\textbf{Perche esiste:} quando il dipendente registra un cliente, questa funzione crea il record.

\textbf{Spiegazione semplice:} e come compilare una nuova scheda cliente e metterla nell'archivio.

\paragraph{getClienti()}
\textbf{Cosa fa:} legge tutti i clienti presenti nel database.

\textbf{Perche esiste:} serve per riempire tabelle e combo box.

\textbf{Spiegazione semplice:} prende tutte le schede e le mette in una lista che Java puo usare.

\paragraph{getClienteByEmail(email)}
\textbf{Cosa fa:} cerca un cliente tramite email.

\textbf{Perche esiste:} e utile per verificare se un cliente esiste gia oppure per recuperarlo rapidamente.

\textbf{Spiegazione semplice:} e come cercare una persona usando la sua email invece del nome.

\paragraph{autenticaCliente(email, password)}
\textbf{Cosa fa:} verifica se email e password coincidono con un cliente reale.

\textbf{Perche esiste:} il login cliente dipende da questa verifica.

\textbf{Spiegazione semplice:} e il controllo all'ingresso: se i dati sono giusti, la porta si apre.

\paragraph{aggiornaCliente(id, nome, tipo, email, telefono)}
\textbf{Cosa fa:} modifica i dati del cliente senza cambiare la password.

\textbf{Perche esiste:} spesso il cliente deve solo aggiornare contatti o nome, non la password.

\textbf{Spiegazione semplice:} e come correggere una scheda gia esistente.

\paragraph{aggiornaClienteConPassword(...)}
\textbf{Cosa fa:} modifica i dati e anche la password.

\textbf{Perche esiste:} serve quando il dipendente vuole cambiare tutto in un colpo solo.

\textbf{Spiegazione semplice:} come aggiornare la scheda completa invece di una sola parte.

\paragraph{eliminaCliente(id)}
\textbf{Cosa fa:} cancella il cliente dal database.

\textbf{Perche esiste:} a volte un record deve essere rimosso per errore, duplicato o richiesta amministrativa.

\textbf{Spiegazione semplice:} toglie la scheda dall'archivio.

\textbf{Chiamate comuni usate nel DAO e perche:}
\begin{itemize}
  \item \texttt{prepareStatement}: mette sicurezza e ordine nelle query.
  \item \texttt{setString/setInt}: inserisce i valori reali al posto dei segnaposto \texttt{?}.
  \item \texttt{executeQuery}: quando si legge.
  \item \texttt{executeUpdate}: quando si scrive o si cancella.
  \item \texttt{ResultSet}: contiene le righe lette dal database.
\end{itemize}

\subsubsection*{DipendenteDAO.java}
\textbf{Ruolo nel progetto:} gestisce il login e la lettura dei dipendenti.

\textbf{Perche serve davvero:} la logica di autenticazione deve essere separata dalla finestra di login.

\textbf{Spiegazione semplice:} e il controllo documenti del lavoratore.

\paragraph{getDipendenteById(id)}
\textbf{Cosa fa:} cerca un dipendente nel database usando il suo numero identificativo.

\textbf{Perche esiste:} il sistema interno usa l'ID invece dell'email.

\textbf{Spiegazione semplice:} e come cercare una persona tramite numero tessera.

\paragraph{autenticaDipendente(id, password)}
\textbf{Cosa fa:} confronta password inserita e password salvata.

\textbf{Perche esiste:} senza questa verifica, chiunque potrebbe entrare nel gestionale.

\textbf{Spiegazione semplice:} se il tesserino e la parola segreta coincidono, il dipendente entra.

\subsubsection*{VeicoloDAO.java}
\textbf{Ruolo nel progetto:} gestisce i mezzi, il loro stato e la disponibilita.

\textbf{Perche serve davvero:} tutte le schermate che mostrano veicoli dipendono da questa classe.

\textbf{Spiegazione semplice:} e il magazzino digitale dei veicoli.

\paragraph{inserisciVeicolo(targa, marca, modello, tipo, prezzo)}
\textbf{Cosa fa:} inserisce un nuovo veicolo.

\textbf{Perche esiste:} il dipendente deve aggiungere nuovi mezzi al catalogo.

\textbf{Spiegazione semplice:} mette una nuova auto o moto nello scaffale del concessionario.

\paragraph{getStatoVeicolo(targa)}
\textbf{Cosa fa:} chiede al database se il veicolo risulta venduto, temporaneamente bloccato o libero.

\textbf{Perche esiste:} lo stato non va inventato dalla grafica, ma calcolato secondo regole certe.

\textbf{Spiegazione semplice:} domanda al database: ``questo mezzo si puo ancora usare o no?''.

\paragraph{getVeicoliDisponibili()}
\textbf{Cosa fa:} restituisce solo i veicoli prenotabili.

\textbf{Perche esiste:} in vendita e test drive bisogna proporre solo mezzi disponibili.

\textbf{Spiegazione semplice:} mostra solo quello che si puo davvero scegliere.

\paragraph{aggiornaVeicolo(targa, marca, modello, tipo, prezzo)}
\textbf{Cosa fa:} modifica i dati del mezzo.

\textbf{Perche esiste:} prezzo o descrizione possono cambiare.

\textbf{Spiegazione semplice:} corregge la scheda del veicolo.

\paragraph{eliminaVeicolo(targa)}
\textbf{Cosa fa:} elimina un veicolo.

\textbf{Perche esiste:} serve quando un mezzo viene rimosso dal catalogo.

\textbf{Spiegazione semplice:} toglie il mezzo dal magazzino.

\paragraph{getAllVeicoli()}
\textbf{Cosa fa:} legge tutti i veicoli e calcola lo stato.

\textbf{Perche esiste:} il catalogo completo deve mostrare tutto, non solo i mezzi disponibili.

\textbf{Spiegazione semplice:} e la lista completa di tutto il parco veicoli.

\subsubsection*{TestDriveDAO.java}
\textbf{Ruolo nel progetto:} salva le richieste di test drive.

\textbf{Perche serve davvero:} il test drive e una vera operazione del sistema, non solo un messaggio a video.

\textbf{Spiegazione semplice:} scrive nel database che un cliente vuole provare un mezzo in una certa data.

\paragraph{richiediTestDrive(clienteId, targa, data)}
\textbf{Cosa fa:} inserisce una riga nella tabella TestDrive.

\textbf{Perche esiste:} la prenotazione va memorizzata e controllata.

\textbf{Spiegazione semplice:} e come segnare su un registro che una prova e stata fissata.

\subsubsection*{VenditaDAO.java}
\textbf{Ruolo nel progetto:} salva una vendita ufficiale.

\textbf{Perche serve davvero:} la vendita e l'evento finale piu importante di un concessionario.

\textbf{Spiegazione semplice:} quando si conclude un acquisto, questa classe lascia la traccia nel database.

\paragraph{effettuaVendita(targa, clienteId, dipendenteId, prezzoFinale)}
\textbf{Cosa fa:} crea il record di vendita con la data corrente.

\textbf{Perche esiste:} la vendita deve contenere chi ha comprato, chi ha venduto e a che prezzo.

\textbf{Spiegazione semplice:} scrive il contratto finale in forma digitale.

\subsection{5. Utility della GUI}

\subsubsection*{UiSupport.java}
\textbf{Ruolo nel progetto:} fornisce funzioni visive e per errori comuni.

\textbf{Perche serve davvero:} la stessa grafica e lo stesso modo di mostrare errori devono essere coerenti in tutto il programma.

\textbf{Spiegazione semplice:} e una piccola scatola degli attrezzi da usare ovunque.

\paragraph{wrapTitled(panel, title)}
\textbf{Cosa fa:} mette un bordo con titolo al pannello.

\textbf{Perche esiste:} aiuta a separare visivamente le aree della finestra.

\textbf{Spiegazione semplice:} e come mettere un'etichetta sopra una scatola.

\paragraph{showErr(parent, ex)}
\textbf{Cosa fa:} mostra un popup di errore.

\textbf{Perche esiste:} l'utente deve capire che qualcosa e andato storto senza leggere dettagli tecnici troppo difficili.

\textbf{Spiegazione semplice:} fa comparire una finestra che dice che c'e stato un problema.

\paragraph{toUserMessage(ex)}
\textbf{Cosa fa:} converte alcuni errori tecnici in messaggi piu semplici.

\textbf{Perche esiste:} il database puo restituire messaggi poco chiari; questa funzione li traduce in modo comprensibile.

\textbf{Spiegazione semplice:} trasforma il linguaggio ``da tecnico'' in linguaggio normale.

\subsubsection*{UiModels.java}
\textbf{Ruolo nel progetto:} aiuta le combo box a mostrare testi leggibili.

\textbf{Perche serve davvero:} se una combo mostra solo numeri, l'utente si confonde facilmente.

\textbf{Spiegazione semplice:} trasforma dati grezzi in etichette facili da leggere.

\paragraph{ClienteItem}
\textbf{Cosa fa:} conserva id e nome.

\textbf{Perche esiste:} la combo del cliente deve mostrare un testo utile, non solo l'id.

\textbf{Spiegazione semplice:} mostra qualcosa come ``Mario Rossi (3)''.

\paragraph{VeicoloItem}
\textbf{Cosa fa:} conserva targa, modello e prezzo.

\textbf{Perche esiste:} in una lista di scelta e piu utile vedere anche il modello del mezzo.

\textbf{Spiegazione semplice:} e come leggere una targhetta con il nome del mezzo.

\paragraph{ClienteLogin}
\textbf{Cosa fa:} conserva un cliente e la password usata nel login.

\textbf{Perche esiste:} e un piccolo contenitore utile nei flussi di autenticazione.

\textbf{Spiegazione semplice:} mette insieme chi prova ad entrare e la chiave che usa.

\subsubsection*{InputValidator.java}
\textbf{Ruolo nel progetto:} controlla la correttezza degli input prima del salvataggio.

\textbf{Perche serve davvero:} senza validazione, il database riceverebbe dati sbagliati e l'app diventerebbe meno affidabile.

\textbf{Spiegazione semplice:} e il professore severo che controlla se quello che scrivi e scritto bene.

\paragraph{isValidEmail(email)}
\textbf{Cosa fa:} controlla se l'email ha una forma accettabile.

\textbf{Perche esiste:} evita errori evidenti come email senza @ o con spazi.

\textbf{Spiegazione semplice:} verifica che l'indirizzo sembri davvero una mail.

\paragraph{isValidPhone(phone)}
\textbf{Cosa fa:} controlla il numero di telefono.

\textbf{Perche esiste:} serve un formato consistente.

\textbf{Spiegazione semplice:} controlla che il numero sia fatto come previsto.

\paragraph{normalizePlate(plate)}
\textbf{Cosa fa:} toglie spazi esterni e converte in maiuscolo.

\textbf{Perche esiste:} le targhe devono essere confrontate in modo uniforme.

\textbf{Spiegazione semplice:} rende tutte le targhe scritte nello stesso modo.

\paragraph{isValidPlate(plate)}
\textbf{Cosa fa:} verifica che la targa abbia caratteri ammessi.

\textbf{Perche esiste:} evita lettere strane, simboli o spazi sbagliati.

\textbf{Spiegazione semplice:} controlla che la targa assomigli davvero a una targa.

\paragraph{parseDateOrThrow(dateInput)}
\textbf{Cosa fa:} trasforma una stringa in data oppure lancia un errore chiaro.

\textbf{Perche esiste:} l'utente scrive date come testo, ma il programma ha bisogno di un oggetto data.

\textbf{Spiegazione semplice:} traduce una data scritta a mano in una data che Java capisce.

\paragraph{requireTodayOrFuture(date)}
\textbf{Cosa fa:} impedisce di usare date passate.

\textbf{Perche esiste:} un test drive non puo essere prenotato in un giorno gia passato.

\textbf{Spiegazione semplice:} evita di viaggiare nel tempo all'indietro.

\subsection{6. Finestra principale e navigazione}

\subsubsection*{ConcessionarioFrame.java}
\textbf{Ruolo nel progetto:} coordina login, logout, viste dipendente e viste cliente.

\textbf{Perche serve davvero:} in un gestionale grafico serve un unico posto che sappia cambiare schermata in base all'utente.

\textbf{Spiegazione semplice:} e il cervello della finestra principale.

\paragraph{ConcessionarioFrame()}
\textbf{Cosa fa:} prepara finestra, header e area login.

\textbf{Perche esiste:} cosi la struttura base viene costruita una sola volta.

\textbf{Spiegazione semplice:} costruisce la casa prima che qualcuno ci entri.

\paragraph{createHeader()}
\textbf{Cosa fa:} crea la barra superiore con titolo e logout.

\textbf{Perche esiste:} l'utente deve sempre capire dove si trova e chi e loggato.

\textbf{Spiegazione semplice:} e il cartello sopra la porta.

\paragraph{createLoginPanel()}
\textbf{Cosa fa:} crea il modulo di accesso con ruolo, credenziali e note.

\textbf{Perche esiste:} e il punto in cui l'utente entra nell'app.

\textbf{Spiegazione semplice:} e il modulo d'ingresso.

\paragraph{buildDipendenteArea(id)}
\textbf{Cosa fa:} crea i tab del dipendente.

\textbf{Perche esiste:} il personale interno ha piu funzioni da gestire.

\textbf{Spiegazione semplice:} apre il pannello di lavoro del dipendente.

\paragraph{buildClienteArea(cliente)}
\textbf{Cosa fa:} crea i tab del cliente.

\textbf{Perche esiste:} un cliente non deve vedere strumenti amministrativi inutili.

\textbf{Spiegazione semplice:} mostra la parte dell'app che il cliente puo usare.

\paragraph{switchToLogin()}
\textbf{Cosa fa:} torna al login.

\textbf{Perche esiste:} dopo il logout bisogna ripulire la sessione.

\textbf{Spiegazione semplice:} riporta l'app alla porta d'ingresso.

\paragraph{switchToDipendente(dipendente)}
\textbf{Cosa fa:} mostra l'area dipendente e aggiorna la sessione.

\textbf{Perche esiste:} dopo il login bisogna portare l'utente nella zona corretta.

\textbf{Spiegazione semplice:} accompagna il lavoratore nel suo ufficio digitale.

\paragraph{switchToCliente(cliente)}
\textbf{Cosa fa:} mostra l'area cliente e aggiorna la sessione.

\textbf{Perche esiste:} il cliente deve vedere solo le sue funzioni.

\textbf{Spiegazione semplice:} porta il cliente nella sua area riservata.

\paragraph{handleClienteLogin(email,password)}
\textbf{Cosa fa:} controlla i campi e prova l'autenticazione cliente.

\textbf{Perche esiste:} cosi si valida l'input prima di parlare con il database.

\textbf{Spiegazione semplice:} prima guarda se i dati hanno senso, poi verifica se l'utente esiste.

\paragraph{handleDipendenteLogin(idText,password)}
\textbf{Cosa fa:} controlla i campi e prova l'autenticazione dipendente.

\textbf{Perche esiste:} l'ID deve essere numerico e la password non vuota.

\textbf{Spiegazione semplice:} e il controllo ingresso per chi lavora nel concessionario.

\subsection{7. Pannelli cliente}

\subsubsection*{CatalogoClientePanel.java}
\textbf{Cosa fa:} mostra il catalogo completo.

\textbf{Perche esiste:} il cliente deve poter guardare cosa offre il concessionario.

\textbf{Spiegazione semplice:} e una vetrina digitale.

\textbf{Comportamento interno spiegato:}
\begin{itemize}
  \item crea una tabella con colonne leggibili,
  \item chiama \texttt{VeicoloDAO.getAllVeicoli()},
  \item ricarica i dati quando si preme refresh.
\end{itemize}

\subsubsection*{ProfiloClientePanel.java}
\textbf{Cosa fa:} permette di leggere e modificare il profilo del cliente.

\textbf{Perche esiste:} le informazioni di contatto possono cambiare.

\textbf{Spiegazione semplice:} e la pagina dei tuoi dati personali.

\textbf{Passo per passo:}
\begin{enumerate}
  \item apre il profilo dal database,
  \item mostra i campi,
  \item controlla email e telefono,
  \item salva le modifiche.
\end{enumerate}

\subsubsection*{TestDriveClientePanel.java}
\textbf{Cosa fa:} prenota un test drive per il cliente loggato.

\textbf{Perche esiste:} il test drive e una funzione tipica del cliente.

\textbf{Spiegazione semplice:} scegli un veicolo, scegli una data, chiedi di provarlo.

\textbf{Perche servono le validazioni:}
senza controllo data si potrebbero inserire richieste sbagliate o impossibili.

\subsubsection*{OperazioniClientePanel.java}
\textbf{Cosa fa:} mostra le vendite e i test drive del cliente.

\textbf{Perche esiste:} l'utente deve poter rivedere la propria cronologia.

\textbf{Spiegazione semplice:} e un archivio personale delle cose gia fatte.

\subsection{8. Pannelli dipendente}

\subsubsection*{ClientiPanel.java}
\textbf{Cosa fa:} gestisce i clienti in modo completo.

\textbf{Perche esiste:} i dipendenti devono poter inserire, modificare e cancellare clienti.

\textbf{Spiegazione semplice:} e il foglio di lavoro dei clienti.

\textbf{Comportamento da studiare:}
\begin{itemize}
  \item la tabella viene popolata da DAO,
  \item la selezione di una riga riempie i campi,
  \item il salvataggio controlla i dati,
  \item l'eliminazione chiede conferma.
\end{itemize}

\subsubsection*{VeicoliPanel.java}
\textbf{Cosa fa:} gestisce il parco veicoli.

\textbf{Perche esiste:} senza questo pannello il concessionario non potrebbe amministrare auto e moto.

\textbf{Spiegazione semplice:} e il cruscotto di tutto il magazzino.

\textbf{Passi importanti:}
\begin{enumerate}
  \item controlla la targa,
  \item controlla la marca,
  \item se la marca manca chiede la nazione,
  \item salva o aggiorna il mezzo,
  \item rinfresca la tabella.
\end{enumerate}

\subsubsection*{VenditaPanel.java}
\textbf{Cosa fa:} registra una vendita con cliente e veicolo scelti.

\textbf{Perche esiste:} e il punto operativo dove una trattativa diventa vendita vera.

\textbf{Spiegazione semplice:} come chiudere un affare e scriverlo nel libro contabile.

\subsubsection*{TestDrivePanel.java}
\textbf{Cosa fa:} fa prenotare test drive anche al personale.

\textbf{Perche esiste:} a volte la prenotazione viene inserita dal dipendente per conto del cliente.

\textbf{Spiegazione semplice:} e la versione interna del pannello test drive.

\subsubsection*{OrdiniPanel.java}
\textbf{Cosa fa:} mostra tutte le vendite e i test drive e permette cancellazioni.

\textbf{Perche esiste:} l'amministrazione deve controllare lo storico completo.

\textbf{Spiegazione semplice:} e la cronologia generale del concessionario.

\subsection{9. SQL, schema, trigger e regole}

\subsubsection*{Tabelle del database}
\textbf{Cosa fanno:} conservano i dati permanenti.

\textbf{Perche esistono:} senza tabelle non ci sarebbe memoria stabile delle operazioni.

\textbf{Spiegazione semplice:} sono gli archivi del concessionario.

\paragraph{Marca}
contiene nome e nazione della marca.

\paragraph{Cliente}
contiene i clienti, il tipo e la password.

\paragraph{Dipendente}
contiene i dipendenti e la password.

\paragraph{Veicolo}
contiene i mezzi e la loro marca.

\paragraph{Vendita}
contiene le vendite gia effettuate.

\paragraph{TestDrive}
contiene le prenotazioni dei test drive.

\subsubsection*{Funzione get\_stato\_veicolo}
\textbf{Cosa fa:} decide lo stato del mezzo.

\textbf{Perche esiste:} la disponibilita dipende da regole e storici, non solo dal dato anagrafico.

\textbf{Spiegazione semplice:} il database controlla se il mezzo e gia stato venduto o impegnato oggi.

\subsubsection*{Trigger di vendita}
\textbf{Cosa fa:} impedisce vendite doppie dello stesso veicolo.

\textbf{Perche esiste:} un veicolo non puo essere venduto due volte.

\textbf{Spiegazione semplice:} e un blocco automatico di sicurezza.

\subsubsection*{Trigger test drive}
\textbf{Cosa fa:} limita il numero mensile di test drive.

\textbf{Perche esiste:} protegge la regola di business del concessionario.

\textbf{Spiegazione semplice:} evita che un cliente prenoti troppe prove nello stesso mese.

\subsection{10. Lettura completa del flusso applicativo}

\textbf{Percorso del cliente:}
\begin{enumerate}
  \item entra dalla finestra di login,
  \item inserisce email e password,
  \item il programma controlla i dati,
  \item il cliente viene portato nella propria area,
  \item puo vedere catalogo, profilo, test drive e storico.
\end{enumerate}

\textbf{Percorso del dipendente:}
\begin{enumerate}
  \item entra dalla finestra di login,
  \item inserisce id e password,
  \item il programma controlla i dati,
  \item il dipendente viene portato nell'area gestionale,
  \item puo lavorare su clienti, veicoli, vendite, test drive e ordini.
\end{enumerate}

\textbf{Perche questa distinzione e utile:}
mantiene il programma semplice da capire e separa bene i permessi.

\subsection{11. Concetti da ricordare durante lo studio}
\begin{itemize}
  \item Una classe non e un oggetto: la classe e il modello, l'oggetto e il risultato concreto.
  \item Un metodo non salva dati da solo: svolge una azione precisa.
  \item Il DAO non mostra la grafica: parla con il database.
  \item La GUI non dovrebbe contenere SQL complessa se non strettamente necessario.
  \item Il DB non e solo memoria: e anche controllo delle regole.
\end{itemize}

\subsection{12. Lettura didattica finale per chi non sa Java}
Se devi spiegare questo progetto a voce, puoi farlo cosi:
\begin{enumerate}
  \item Il programma si apre da una piccola classe iniziale.
  \item Questa classe lancia la finestra principale.
  \item La finestra mostra un login.
  \item In base al ruolo, l'utente entra in una parte diversa dell'app.
  \item Ogni pulsante chiama un metodo.
  \item I metodi o controllano i dati o parlano con il database.
  \item Il database applica anche regole proprie, come limiti e blocchi.
\end{enumerate}

\textbf{Perche questa sequenza e importante:} capire il flusso generale rende piu facile capire ogni singolo file.

\subsection{13. Lettura finale delle motivazioni progettuali}
\begin{itemize}
  \item Si usa Java per avere un'app desktop strutturata e portabile.
  \item Si usa Swing per costruire finestre e componenti grafici.
  \item Si usa JDBC per collegarsi a PostgreSQL.
  \item Si usano DAO per separare SQL e interfaccia.
  \item Si usano model per rappresentare i dati.
  \item Si usano utility per non ripetere codice.
  \item Si usano trigger e funzioni SQL per proteggere le regole fondamentali.
\end{itemize}

\textbf{Spiegazione semplice:} il progetto e organizzato come una piccola azienda digitale: c'e chi mostra la finestra, chi fa il lavoro tecnico, chi conserva i dati e chi controlla che le regole siano rispettate.

\newpage
\section{Estensione avanzata richiesta: spiegazione molto piu approfondita}
Questa sezione e stata aggiunta per rendere la documentazione piu lunga, piu tecnica e piu didattica.
L'obiettivo e spiegare:
\begin{itemize}
  \item cos'e il progetto in termini architetturali e operativi;
  \item cos'e Swing in modo piu serio (non solo definizione breve);
  \item cos'e JDBC con focus pratico su questa codebase;
  \item perche vengono usate API e simboli specifici (con riferimenti file/riga);
  \item tutti gli import usati nei file Java, indicando il motivo del loro uso.
\end{itemize}

\subsection{Che cos'e davvero questo progetto (livello architetturale)}
Il progetto e un gestionale desktop Java per concessionario auto/moto con due ruoli applicativi:
\texttt{cliente} e \texttt{dipendente}.

Dal punto di vista architetturale non e un monolite disordinato: adotta una separazione chiara in livelli.
\begin{enumerate}
  \item \textbf{UI Layer (Swing)}: contiene finestra principale e pannelli; gestisce input utente, tabelle, pulsanti, feedback.
  \item \textbf{Application/Domain Layer leggero}: contiene le decisioni operative immediate (validazioni, scelta del flusso login, aggiornamento viste).
  \item \textbf{Data Access Layer (DAO + JDBC)}: isola le query SQL, mappa risultati in oggetti Java, scrive/legge il database.
  \item \textbf{Database Layer (PostgreSQL + trigger/funzioni)}: custodisce integrita, regole vincolanti e persistenza dei dati.
\end{enumerate}

Il valore didattico principale e che mostra in un progetto piccolo/medio una divisione pulita delle responsabilita:
la UI non e obbligata a conoscere i dettagli SQL e i DAO non devono conoscere la logica grafica della finestra.

\subsection{Che cos'e Swing, spiegato bene e collegato al codice}
Swing e il toolkit grafico classico di Java per applicazioni desktop.
Nel progetto viene usato in modalita tipica enterprise didattica:
\begin{itemize}
  \item una \texttt{JFrame} principale che coordina;
  \item tanti \texttt{JPanel} specializzati;
  \item layout manager per comporre schermate flessibili;
  \item listener evento per reagire alle azioni utente;
  \item modelli tabellari per separare dati e presentazione.
\end{itemize}

\textbf{Perche Swing e adatto qui:}
\begin{itemize}
  \item non richiede browser o server web;
  \item permette form complessi, tabelle e workflow amministrativi;
  \item integra bene JDBC in applicazioni standalone;
  \item e coerente con un progetto universitario di POO + DB.
\end{itemize}

\textbf{Componenti Swing/AWT chiave con uso concreto nel progetto:}
\begin{longtable}{p{0.19\textwidth} p{0.31\textwidth} p{0.42\textwidth}}
\toprule
\textbf{Componente} & \textbf{Dove (file:riga)} & \textbf{Perche viene usato} \\
\midrule
\endhead
\texttt{JFrame} & \texttt{ui/ConcessionarioFrame.java} (classe principale) & Finestra contenitore dell'intera applicazione.
\\
\texttt{JPanel} & \texttt{ui/ConcessionarioFrame.java:28,43,54,80,124,160,173} e in tutti i pannelli sotto \texttt{ui/panels/*} & Unita base di composizione: ogni area funzionale e un pannello riutilizzabile.
\\
\texttt{CardLayout} & \texttt{ui/ConcessionarioFrame.java:27} & Cambio vista login/cliente/dipendente senza aprire nuove finestre.
\\
\texttt{JTabbedPane} & \texttt{ui/ConcessionarioFrame.java:153,167}; \texttt{ui/panels/customer/OperazioniClientePanel.java:20}; \texttt{ui/panels/employee/OrdiniPanel.java:21} & Organizza funzionalita in schede tematiche mantenendo la stessa finestra.
\\
\texttt{FlowLayout} & \texttt{ui/ConcessionarioFrame.java:69,124}; pannelli operativi (es. \texttt{employee/VeicoliPanel.java:54}, \texttt{customer/TestDriveClientePanel.java:34}) & Disposizione lineare dei pulsanti nelle barre azioni.
\\
\texttt{GridBagLayout} & \texttt{ui/ConcessionarioFrame.java:80,83}; vari pannelli form & Layout a griglia flessibile per form con etichette/campi.
\\
\texttt{JButton} & \texttt{ui/ConcessionarioFrame.java:30,123}; in tutti i pannelli operativi & Punto di attivazione esplicito dei casi d'uso (salva, aggiorna, elimina, conferma).
\\
\texttt{addActionListener} & \texttt{ui/ConcessionarioFrame.java:67,103,134}; numerose righe nei pannelli (es. \texttt{employee/VeicoliPanel.java:102,118,136,154,155}) & Collega l'evento click alla logica applicativa.
\\
\bottomrule
\end{longtable}

\subsection{Che cos'e JDBC, spiegato bene e collegato al codice}
JDBC e l'API standard Java per dialogare con database relazionali.
Nel progetto si vede il ciclo completo:
\begin{enumerate}
  \item apertura connessione;
  \item preparazione query (spesso parametrizzate);
  \item esecuzione;
  \item lettura del risultato;
  \item chiusura risorse con \texttt{try-with-resources}.
\end{enumerate}

\textbf{Perche e importante didatticamente:}
\begin{itemize}
  \item fa vedere separazione tra SQL e UI;
  \item riduce memory leak grazie alla chiusura automatica;
  \item rende esplicito il mapping \texttt{ResultSet -> oggetti model}.
\end{itemize}

\subsection{API/simboli richiesti: dove sono usati e perche}
In tabella seguente i duplicati identici sono compressi.
Quando la stessa API e usata con scopi diversi, i casi sono separati esplicitamente.

\begin{longtable}{p{0.18\textwidth} p{0.33\textwidth} p{0.41\textwidth}}
\toprule
\textbf{Simbolo/API} & \textbf{Riferimenti (file:riga)} & \textbf{Motivazione d'uso nel progetto} \\
\midrule
\endhead
\texttt{getConnection()} & \texttt{DBManager.java:12}; implementazione reale \texttt{DBManager.java:13} con \texttt{DriverManager.getConnection(...)} & Centralizza l'apertura connessioni DB e impedisce duplicazione di URL/credenziali.
\\
\texttt{DBManager.getConnection} & DAO: \texttt{dao/ClienteDAO.java:17,31,49,69,93,106,120}; \texttt{dao/DipendenteDAO.java:17}; \texttt{dao/VeicoloDAO.java:17,30,45,68,81,91}; \texttt{dao/TestDriveDAO.java:12}; \texttt{dao/VenditaDAO.java:12}; UI con query diretta: \texttt{employee/OrdiniPanel.java:55,105,135}, \texttt{employee/VeicoliPanel.java:197,233,245}, \texttt{customer/ProfiloClientePanel.java:54,77}, \texttt{customer/OperazioniClientePanel.java:55} & Accesso uniforme al database in tutti i casi d'uso: CRUD clienti/veicoli, autenticazione, storico, vendite e test drive.
\\
\texttt{SQLException} & Import e firma in \texttt{DBManager.java:5,12}; in tutti i DAO (es. \texttt{ClienteDAO.java:9}); nei metodi panel dove servono query & Propagazione formale degli errori SQL per gestirli in UI con messaggi utente coerenti.
\\
\texttt{createStatement()} & \texttt{dao/ClienteDAO.java:32}; \texttt{dao/VeicoloDAO.java:46,92}; \texttt{employee/OrdiniPanel.java:55} & Usato quando non servono parametri in input (SELECT di lettura/reportistica).
\\
\texttt{Statement} & Import in \texttt{dao/ClienteDAO.java:10}, \texttt{dao/VeicoloDAO.java:10}, \texttt{employee/OrdiniPanel.java:11} & Permette query statiche semplici per caricamento liste o viste aggregate.
\\
\texttt{getString(...)} su \texttt{ResultSet} & Esempi: \texttt{dao/ClienteDAO.java:37-40,56-59,76,82-85}; \texttt{dao/DipendenteDAO.java:24-26}; \texttt{dao/VeicoloDAO.java:35,49,54-56,95,99-101}; \texttt{employee/OrdiniPanel.java:60-63,72-74}; \texttt{customer/ProfiloClientePanel.java:61-64}; \texttt{customer/OperazioniClientePanel.java:64-65,67,77-78} & Estrae colonne testuali (o date rese stringa in query) e alimenta model Java o tabelle Swing.
\\
\texttt{FlowLayout} & \texttt{ui/ConcessionarioFrame.java:69,124}; \texttt{customer/CatalogoClientePanel.java:18}; \texttt{customer/TestDriveClientePanel.java:34}; \texttt{customer/ProfiloClientePanel.java:43}; \texttt{employee/ClientiPanel.java} (toolbar), \texttt{employee/VeicoliPanel.java:54}, \texttt{employee/TestDrivePanel.java:40}, \texttt{employee/OrdiniPanel.java:44} & Layout semplice per barre pulsanti allineate e facilmente leggibili.
\\
\texttt{JPanel} & Dichiarato/istanziato in tutta la UI, es. \texttt{ui/ConcessionarioFrame.java:28,43,54,80,124,160,173} & Elemento base per suddividere UI in blocchi indipendenti e riusabili.
\\
\texttt{actions} (variabile pannello pulsanti) & \texttt{ui/ConcessionarioFrame.java:124-128}; \texttt{customer/TestDriveClientePanel.java:34-39}; \texttt{customer/ProfiloClientePanel.java:43-47}; \texttt{employee/VeicoliPanel.java:54-63}; \texttt{employee/TestDrivePanel.java:40-45} & Convenzione locale per raggruppare pulsanti di comando in una toolbar coerente.
\\
\texttt{JButton} & \texttt{ui/ConcessionarioFrame.java:30,123}; uso diffuso nei pannelli (refresh/salva/elimina/conferma) & Trigger esplicito delle azioni utente.
\\
\texttt{loginBtn} & \texttt{ui/ConcessionarioFrame.java:123,126,134} & Pulsante di accesso: definizione, inserimento nel pannello azioni, bind al listener di login.
\\
\texttt{UiSupport} & Definizione \texttt{ui/common/UiSupport.java:6}; uso in \texttt{ui/ConcessionarioFrame.java:7,83,143} e in quasi tutti i pannelli & Standardizza look dei blocchi titolati e gestione errori in popup.
\\
\texttt{handleDipendenteLogin} & Definizione \texttt{ui/ConcessionarioFrame.java:227}; invocazione nel routing login \texttt{ui/ConcessionarioFrame.java:140} & Incapsula validazione ID/password dipendente + autenticazione DAO + switch vista.
\\
\texttt{equalsIgnoreCase} & \texttt{ui/ConcessionarioFrame.java:105,137} & Confronto robusto del ruolo selezionato (evita problemi maiuscole/minuscole).
\\
\texttt{addActionListener} & \texttt{ui/ConcessionarioFrame.java:67,103,134}; molti esempi nei pannelli cliente/dipendente & Meccanismo evento-azione di Swing per comportamento reattivo della GUI.
\\
\texttt{JTabbedPane} & \texttt{ui/ConcessionarioFrame.java:153,167}; \texttt{customer/OperazioniClientePanel.java:20}; \texttt{employee/OrdiniPanel.java:21} & Gestione di sotto-sezioni in tab senza moltiplicare finestre.
\\
\bottomrule
\end{longtable}

\subsection{Stesso simbolo, motivi diversi: casi esplicitati}
Di seguito gli esempi richiesti in cui lo stesso costrutto appare in contesti differenti.

\begin{longtable}{p{0.22\textwidth} p{0.31\textwidth} p{0.39\textwidth}}
\toprule
\textbf{Costrutto} & \textbf{Contesto} & \textbf{Differenza di scopo} \\
\midrule
\endhead
\texttt{createStatement()} & \texttt{dao/ClienteDAO.java:32}, \texttt{dao/VeicoloDAO.java:46,92} & Caricamento liste standard (clienti/veicoli) senza parametri runtime.
\\
\texttt{createStatement()} & \texttt{employee/OrdiniPanel.java:55} & Query aggregate con JOIN per vista amministrativa (storico ordini globale).
\\
\texttt{JTabbedPane} & \texttt{ui/ConcessionarioFrame.java:153,167} & Navigazione macro tra funzionalita principali per ruolo.
\\
\texttt{JTabbedPane} & \texttt{employee/OrdiniPanel.java:21}, \texttt{customer/OperazioniClientePanel.java:20} & Suddivisione interna di uno stesso pannello in sotto-viste omogenee.
\\
\texttt{addActionListener} & \texttt{ui/ConcessionarioFrame.java:134} su \texttt{loginBtn} & Orchestrazione autenticazione e scelta flusso ruolo.
\\
\texttt{addActionListener} & \texttt{employee/VeicoliPanel.java:102,118,136,154,155} & Comandi CRUD multipli sul catalogo veicoli.
\\
\bottomrule
\end{longtable}

\subsection{Mappa completa import usati, file per file}
In questa sezione ogni file Java viene coperto.
Se un file non ha import, viene indicato esplicitamente.

\subsubsection*{Package \texttt{it.unina.prog}}
\paragraph{\texttt{DBManager.java}}
\begin{itemize}
  \item \texttt{java.sql.Connection} (riga 3): tipo di ritorno del canale DB.
  \item \texttt{java.sql.DriverManager} (riga 4): factory JDBC per apertura connessione.
  \item \texttt{java.sql.SQLException} (riga 5): eccezione tecnica propagata dal livello DB.
\end{itemize}

\paragraph{\texttt{Main.java}}
Nessun import: classe bootstrap minimale che delega a \texttt{MainApp}.

\paragraph{\texttt{MainApp.java}}
\begin{itemize}
  \item \texttt{javax.swing.SwingUtilities} (riga 3): avvio GUI su EDT.
  \item \texttt{javax.swing.UIManager} (riga 4): look and feel nativo.
  \item \texttt{it.unina.prog.ui.ConcessionarioFrame} (riga 6): finestra principale.
\end{itemize}

\subsubsection*{Package \texttt{it.unina.prog.model}}
\paragraph{\texttt{Cliente.java}} Nessun import (POJO puro).
\paragraph{\texttt{Dipendente.java}} Nessun import (POJO puro).
\paragraph{\texttt{Veicolo.java}} Nessun import (POJO puro).

\subsubsection*{Package \texttt{it.unina.prog.dao}}
\paragraph{\texttt{ClienteDAO.java}}
\begin{itemize}
  \item \texttt{it.unina.prog.DBManager} (riga 3): connessione centralizzata.
  \item \texttt{it.unina.prog.model.Cliente} (riga 4): mapping righe DB in oggetti.
  \item \texttt{java.sql.Connection} (riga 6): sessione DB.
  \item \texttt{java.sql.PreparedStatement} (riga 7): query parametrizzate sicure.
  \item \texttt{java.sql.ResultSet} (riga 8): lettura risultati SELECT.
  \item \texttt{java.sql.SQLException} (riga 9): gestione errori SQL.
  \item \texttt{java.sql.Statement} (riga 10): query semplici senza parametri.
  \item \texttt{java.util.ArrayList} (riga 11), \texttt{java.util.List} (riga 12): collezioni di clienti.
\end{itemize}

\paragraph{\texttt{DipendenteDAO.java}}
\begin{itemize}
  \item \texttt{it.unina.prog.DBManager} (riga 3), \texttt{it.unina.prog.model.Dipendente} (riga 4).
  \item \texttt{Connection}, \texttt{PreparedStatement}, \texttt{ResultSet}, \texttt{SQLException} (righe 6--9): pipeline login/lookup dipendente.
\end{itemize}

\paragraph{\texttt{VeicoloDAO.java}}
\begin{itemize}
  \item \texttt{it.unina.prog.DBManager} (riga 3), \texttt{it.unina.prog.model.Veicolo} (riga 4).
  \item \texttt{Connection}, \texttt{PreparedStatement}, \texttt{ResultSet}, \texttt{SQLException}, \texttt{Statement} (righe 6--10): CRUD + query catalogo/stato.
  \item \texttt{ArrayList}, \texttt{List} (righe 11--12): raccolta veicoli per UI.
\end{itemize}

\paragraph{\texttt{TestDriveDAO.java}}
\begin{itemize}
  \item \texttt{it.unina.prog.DBManager} (riga 3).
  \item \texttt{Connection}, \texttt{PreparedStatement}, \texttt{SQLException} (righe 5--7): inserimento prenotazione test drive.
\end{itemize}

\paragraph{\texttt{VenditaDAO.java}}
\begin{itemize}
  \item \texttt{it.unina.prog.DBManager} (riga 3).
  \item \texttt{Connection}, \texttt{PreparedStatement}, \texttt{SQLException} (righe 5--7): registrazione vendita.
\end{itemize}

\subsubsection*{Package \texttt{it.unina.prog.ui.common}}
\paragraph{\texttt{UiSupport.java}}
\begin{itemize}
  \item \texttt{javax.swing.*} (riga 3): popup e bordi titolati.
  \item \texttt{java.awt.*} (riga 4): tipo \texttt{Component} e colori.
\end{itemize}

\subsubsection*{Package \texttt{it.unina.prog.ui.model}}
\paragraph{\texttt{UiModels.java}}
\begin{itemize}
  \item \texttt{it.unina.prog.model.Cliente} (riga 3): wrapper UI per selezione cliente.
\end{itemize}

\subsubsection*{Package \texttt{it.unina.prog.ui.validation}}
\paragraph{\texttt{InputValidator.java}}
\begin{itemize}
  \item \texttt{java.time.LocalDate} (riga 3): parsing/controllo date future.
\end{itemize}

\subsubsection*{Package \texttt{it.unina.prog.ui}}
\paragraph{\texttt{ConcessionarioFrame.java}}
\begin{itemize}
  \item DAO import (righe 3--4): autenticazione per ruolo.
  \item Model import (righe 5--6): dati sessione utente autenticato.
  \item \texttt{UiSupport} (riga 7): coerenza UI e error handling.
  \item Import pannelli customer/employee (righe 8--16): composizione dell'area a tab per i due ruoli.
  \item \texttt{javax.swing.*} (riga 18): componenti grafici principali.
  \item \texttt{java.awt.*} (riga 19): layout manager, dimensioni, colori.
\end{itemize}

\subsubsection*{Package \texttt{it.unina.prog.ui.panels.customer}}
\paragraph{\texttt{CatalogoClientePanel.java}}
\begin{itemize}
  \item \texttt{VeicoloDAO}, \texttt{Veicolo}, \texttt{UiSupport} (righe 3--5): dati catalogo + supporto visuale.
  \item \texttt{javax.swing.*} (riga 7), \texttt{DefaultTableModel} (riga 8), \texttt{java.awt.*} (riga 9): tabella e layout.
\end{itemize}

\paragraph{\texttt{OperazioniClientePanel.java}}
\begin{itemize}
  \item \texttt{DBManager} e \texttt{UiSupport} (righe 3--4): query storico + gestione errori.
  \item \texttt{javax.swing.*}, \texttt{DefaultTableModel}, \texttt{java.awt.*} (righe 6--8): tab, tabella e toolbar.
  \item \texttt{Connection}, \texttt{ResultSet} (righe 9--10): lettura storico vendite/test drive.
\end{itemize}

\paragraph{\texttt{ProfiloClientePanel.java}}
\begin{itemize}
  \item \texttt{DBManager}, \texttt{UiSupport}, \texttt{InputValidator} (righe 3--5): caricamento/salvataggio profilo validato.
  \item \texttt{javax.swing.*}, \texttt{java.awt.*} (righe 7--8): form profilo.
  \item \texttt{Connection}, \texttt{ResultSet} (righe 9--10): fetch dati utente.
\end{itemize}

\paragraph{\texttt{TestDriveClientePanel.java}}
\begin{itemize}
  \item \texttt{TestDriveDAO}, \texttt{VeicoloDAO}, \texttt{Veicolo} (righe 3--5): selezione mezzo e conferma prenotazione.
  \item \texttt{UiSupport}, \texttt{UiModels.VeicoloItem}, \texttt{InputValidator} (righe 6--8): UX, etichette combo, regole data.
  \item \texttt{javax.swing.*}, \texttt{java.awt.*}, \texttt{LocalDate} (righe 10--12): form e gestione data.
\end{itemize}

\subsubsection*{Package \texttt{it.unina.prog.ui.panels.employee}}
\paragraph{\texttt{ClientiPanel.java}}
\begin{itemize}
  \item \texttt{ClienteDAO}, \texttt{Cliente}, \texttt{UiSupport}, \texttt{InputValidator} (righe 3--6): CRUD clienti validato.
  \item \texttt{javax.swing.*}, \texttt{DefaultTableModel}, \texttt{java.awt.*}, \texttt{List} (righe 8--11): tabella, layout e lista dati.
\end{itemize}

\paragraph{\texttt{VeicoliPanel.java}}
\begin{itemize}
  \item \texttt{DBManager}, \texttt{VeicoloDAO}, \texttt{Veicolo}, \texttt{UiSupport}, \texttt{InputValidator} (righe 3--7): CRUD veicoli + controlli su marca/targa.
  \item \texttt{javax.swing.*}, \texttt{DefaultTableModel}, \texttt{java.awt.*} (righe 9--11): UI tabellare/form complesso.
  \item \texttt{Connection}, \texttt{ResultSet} (righe 12--13): query di supporto non incapsulate nel DAO.
  \item \texttt{List} (riga 14): caricamento catalogo in memoria.
\end{itemize}

\paragraph{\texttt{VenditaPanel.java}}
\begin{itemize}
  \item \texttt{ClienteDAO}, \texttt{VeicoloDAO}, \texttt{VenditaDAO} (righe 3--5): alimentazione combo + commit vendita.
  \item \texttt{Cliente}, \texttt{Veicolo} (righe 6--7): entita mostrate/selezionate.
  \item \texttt{UiSupport}, \texttt{UiModels.ClienteItem}, \texttt{UiModels.VeicoloItem} (righe 8--10): UX errori + label leggibili in combo.
  \item \texttt{javax.swing.*}, \texttt{java.awt.*} (righe 12--13): form di vendita.
\end{itemize}

\paragraph{\texttt{TestDrivePanel.java}}
\begin{itemize}
  \item \texttt{ClienteDAO}, \texttt{TestDriveDAO}, \texttt{VeicoloDAO} (righe 3--5): raccolta dati e scrittura prenotazione.
  \item \texttt{Cliente}, \texttt{Veicolo} (righe 6--7): tipi dominio in selezione.
  \item \texttt{UiSupport}, \texttt{UiModels.ClienteItem}, \texttt{UiModels.VeicoloItem}, \texttt{InputValidator} (righe 8--11): usabilita e validazione.
  \item \texttt{javax.swing.*}, \texttt{java.awt.*}, \texttt{LocalDate} (righe 13--15): form e regole temporali.
\end{itemize}

\paragraph{\texttt{OrdiniPanel.java}}
\begin{itemize}
  \item \texttt{DBManager}, \texttt{UiSupport} (righe 3--4): query storico globale + error handling.
  \item \texttt{javax.swing.*}, \texttt{DefaultTableModel}, \texttt{java.awt.*} (righe 6--8): doppia tabella e comandi operativi.
  \item \texttt{Connection}, \texttt{ResultSet}, \texttt{Statement} (righe 9--11): SELECT aggregate e cancellazioni amministrative.
\end{itemize}

\subsection{Note finali sulla deduplicazione richiesta}
I duplicati puramente meccanici non sono stati ripetuti riga per riga.
Sono invece stati mantenuti i duplicati \textbf{semanticamente diversi} (stessa API, scopo differente),
come richiesto: esempio \texttt{createStatement()} in DAO di elenco vs \texttt{createStatement()} in storico ordini con JOIN amministrative.

\subsection{Risultato della revisione}
Con questa estensione:
\begin{itemize}
  \item la documentazione supera significativamente la versione precedente in lunghezza e dettaglio;
  \item ogni simbolo/API chiave richiesto e argomentato con riferimenti file/riga;
  \item gli import sono coperti per tutti i file Java del progetto, inclusi quelli senza import.
\end{itemize}

\end{document}
